% Owner: @vinle0 (track C). Only the owner edits this file; send input to them on #40. Functional regression and the attack tests.

% Issue #10 (C2): riscv-tests and Spike co-simulation in the regression script.
\paragraph{Functional regression}
Every configuration in Table~\ref{tab:configs} must pass the RISC-V ISA tests
and the riscv-tests benchmarks in Verilator, as the baseline does. Because a
delay mechanism can fail by stalling forever rather than by computing a wrong
value, a test that times out counts as a failure. This is crucial for catching
deadlocks such as a transmitter that waits on a taint that is never cleared. 
We also run Chipyard's co-simulation mode, which checks each committed instruction
against the Spike ISA simulator in lockstep, so a divergence is reported at
the first wrong instruction rather than at the end of a test. With every new
switch turned off, the generated Verilog must be identical to the baseline's,
which we check by diffing it; this confirms that our changes leave the
unmodified design untouched and baseline numbers stay comparable.

% Issues #9 (C1), #6 (B1).
\paragraph{Security checks}
The Spectre proofs of concept are the end-to-end security tests which all print
whether they recovered the secret and enable the regression script to check the
result automatically. An attack must leak on the baseline, which shows the
test itself works, and must fail on every protected configuration. Since the
conservative stand-in is strictly stronger than STT, it gives a second reference point 
such that any attack that succeeds against it points to a path that bypasses the issue queue, 
rather than to a flaw in taint tracking.
